\chapter{Beyond VHDL: The Design-Side Additions}
\label{ch:beyond}

\chapabstract{\cref{ch:rosetta} and \cref{ch:rosetta2} were a translation
exercise: every \VHDL{} construct had a \SV{} counterpart, and the job was to
find it. This chapter covers the constructs with no \VHDL{} source to
translate from: interfaces, clocking blocks, the checked \code{always}
variants, packed structures and unions, typed enumerations, parameterised
types, the preprocessor, \code{generate} and \code{bind}. Every one of them
elaborates into hardware or into the hierarchy around it. At the end you will
be able to read a modern \SV{} design, and to write one that is shorter and
safer than the \VHDL{} you would have written instead.}

A large part of \SV{} is Verilog with the sharp edges filed off, and the two
preceding chapters covered that part. What remains is genuinely new, and it
splits cleanly in two.

The first half is design-side, and it is this chapter. The order is
deliberate. Interfaces come first, because they change the shape of a
hierarchy and everything after them can be declared inside one. Then come the
constructs that change the shape of a single module: intent-carrying
\code{always} forms, packed structures, typed enumerations, type parameters,
and the preprocessor you should mostly avoid. \code{generate} and \code{bind}
close the chapter, because \kw{bind} is the one construct here that changes
what you can attach to a module you did not write, and it is the bridge to
everything that follows.

The second half is verification-side: dynamic data types, classes,
constrained randomisation, functional coverage, assertions and dynamic
threads. \VHDL{} has nothing comparable, and it is the reason the industry
verifies in \SV{} even where it designs in \VHDL{}. That half is
\refverificationbook, immediately after this one, because none of it makes sense
until you have seen what it is for.

One question belongs to neither half: what does \VHDL{} still do better? That
is a question about the two languages rather than about \SV{} alone, and
\cref{ch:philosophy} answers it, next to the account of where the two
languages came from.

% ===========================================================================
\section{Interfaces}
\label{sec:beyond:interfaces}
\index{interface}

This is the single biggest structural addition to the language, and the one
that changes how you write hierarchy.

Consider a small memory-mapped bus: request, write enable, address, write
data, byte enables, grant, read valid, read data, error. Nine signals, plus a
clock and a reset. In \VHDL{} you name every one of them at least three
times: once in the entity of the master, once in the entity of the slave, and
once in each port map. Add a byte-enable signal six months later and you edit
four files. Get the direction wrong in one of them and the compiler tells you
nothing until elaboration, if at all.

\SV{} lets you declare the bundle once, as a first-class design object:

\begin{svcode}[caption={A small memory-mapped bus as an interface},label={lst:beyond:iface}]
interface simple_bus #(parameter int AW = 12,
                       parameter int DW = 32)
                     (input logic clk, input logic rst_n);

  logic            req;      // master -> slave
  logic            we;
  logic [AW-1:0]   addr;
  logic [DW-1:0]   wdata;
  logic [DW/8-1:0] be;
  logic            gnt;      // slave  -> master
  logic            rvalid;
  logic [DW-1:0]   rdata;
  logic            err;

  // One write transaction, written once, usable by every master.
  task automatic write(input logic [AW-1:0] a,
                       input logic [DW-1:0] d);
    @(posedge clk);
    req <= 1'b1; we <= 1'b1; addr <= a; wdata <= d; be <= '1;
    do @(posedge clk); while (gnt !== 1'b1);
    req <= 1'b0; we <= 1'b0;
  endtask

  modport master  (input  clk, rst_n, gnt, rvalid, rdata, err,
                   output req, we, addr, wdata, be,
                   import write);

  modport slave   (input  clk, rst_n, req, we, addr, wdata, be,
                   output gnt, rvalid, rdata, err);

  modport monitor (input  clk, rst_n, req, we, addr, wdata, be,
                          gnt, rvalid, rdata, err);
endinterface
\end{svcode}

Three things are happening in \cref{lst:beyond:iface}, and they are worth
separating.

\textbf{The bundle.} The nine \kw{logic} declarations are the signals
themselves. An interface instance is a real object in the hierarchy, like a
module instance: it has a name, it can be parameterised, and it can be found
in the wave viewer at \code{top.bus.addr}.

\textbf{The direction views.} A \kw{modport}\index{modport} is a named view of
the bundle from one participant's point of view. It carries no signals of its
own; it says which of the interface's signals that participant may drive and
which it may only read. The same wires are \code{output} in the
\code{master} view and \code{input} in the \code{slave} view. Declaring the
directions once, in one place, is exactly the property \VHDL{} record ports
lack.

\textbf{The behaviour.} The \code{write} task lives inside the interface and
drives the interface's own signals. Any master that imports it through its
modport can call \code{bus.write(12'h040, 32'hDEAD\_BEEF)} and get a
protocol-correct transaction. The protocol is written once and cannot drift
between testbench and design.

Instantiation collapses accordingly:

\begin{svcode}[caption={Instantiating the interface: one connection instead of fifteen},label={lst:beyond:ifaceinst}]
module ram (simple_bus.slave b);
  // inside, the signals are b.req, b.addr, b.wdata, ...
endmodule

module top;
  logic clk, rst_n;

  simple_bus #(.AW(12), .DW(32)) bus (.clk(clk), .rst_n(rst_n));

  cpu     u_cpu (.b(bus.master));
  ram     u_ram (.b(bus.slave));
  bus_mon u_mon (.b(bus.monitor));
endmodule
\end{svcode}

\begin{figure}[htbp]
  \centering
  \begin{tikzpicture}[
      blk/.style={draw=codeframe,line width=1pt,fill=codebg,rounded corners=2pt,
                  minimum width=20mm,minimum height=22mm,align=center,
                  font=\small\sffamily},
      w/.style={draw=codekey,line width=0.5pt,-{Latex[length=1.3mm]}},
      lbl/.style={font=\scriptsize\sffamily,align=center}]

    % --- left: discrete ports ---------------------------------------------
    \node[blk] (a1) at (0,0)   {cpu};
    \node[blk] (a2) at (3.6,0) {ram};
    \foreach \i in {0,...,7} {
      \draw[w] ($(a1.east)+(0,1.05-0.3*\i)$) -- ($(a2.west)+(0,1.05-0.3*\i)$);
    }
    \node[lbl] at (1.8,-1.75) {15 named ports,\\declared in three places};

    % --- right: one interface ---------------------------------------------
    \node[blk] (b1) at (7.6,0)  {cpu};
    \node[blk] (b2) at (11.2,0) {ram};
    \draw[line width=3pt,draw=accent] (b1.east) -- (b2.west);
    \node[draw=accent,fill=white,line width=0.8pt,rounded corners=2pt,
          font=\scriptsize\ttfamily,inner sep=3pt] at (9.4,0) {simple\_bus};
    \node[lbl,anchor=south] at (8.5,0.35) {.master};
    \node[lbl,anchor=south] at (10.4,0.35) {.slave};
    \node[lbl] at (9.4,-1.75) {one connection,\\directions declared once};
  \end{tikzpicture}
  \caption{The same connection with discrete ports and with an interface. The
  modport names on the right are the only place a direction is written.}
  \label{fig:beyond:iface}
\end{figure}

\Cref{fig:beyond:iface} makes the structural point. What it does not show is
the maintenance point: adding a signal to the bus means editing one file, and
every module that does not use the new signal needs no change at all.

\subsection{Parameterised interfaces}

An interface takes parameters exactly as a module does, and the parameters
propagate into every module that receives it. Inside \code{ram}, the width of
\code{b.addr} is whatever \code{AW} was at the instance. You can read the
parameter back through the port:

\begin{svcode}
module ram (simple_bus.slave b);
  localparam int DEPTH = 1 << $bits(b.addr);
  logic [$bits(b.wdata)-1:0] mem [DEPTH];
endmodule
\end{svcode}

This is the closest \SV{} gets to a \VHDL{} generic package: the interface
carries both the types and the objects, and the module adapts to them without
declaring a single generic of its own.

\subsection{Tasks in interfaces, and virtual interfaces}

A modport can \kw{import} a task or a function, and it can also \kw{export}
one, meaning that the interface declares a prototype and the connected module
supplies the body. Importing is the common case and the one to learn first.

Interfaces become far more powerful once classes enter the picture
(\refverificationbook). A class-based testbench holds a \emph{virtual
interface}\index{virtual interface}: a handle to an interface instance, stored
in a variable, assigned at run time. That is what lets a testbench object
which lives in no hierarchy at all reach down and wiggle a wire.
\refverificationbook builds a complete testbench on the mechanism. The declaration
looks like this, and it is the only bridge between the class world and the
module world:

\begin{svcode}
class bus_driver;
  virtual simple_bus.master vif;   // a handle, not an instance
  function new(virtual simple_bus.master v); vif = v; endfunction
  task run(); vif.write(12'h10, 32'h1); endtask
endclass
\end{svcode}

\begin{notebox}{Interfaces are not entities}
An interface is not a lightweight module. It may contain \code{initial} and
\code{always} blocks, assertions, covergroups, tasks, functions, parameters,
types, and even other interface instances. Some verification IP is delivered
almost entirely as an interface. Keep design interfaces boring: signals,
modports, and at most a clocking block.
\end{notebox}

\subsection{What synthesises, and what does not}

Interfaces are synthesisable, and every major synthesis tool in current use
accepts them. They are still avoided in some flows: certain sign-off,
formal-equivalence and legacy scripting flows expect flat port lists, and some
IP-delivery rules forbid them outright because the customer's tool chain may
be older than yours. The pragmatic answer is to keep to a safe subset:

\begin{itemize}
  \item interface instances that are static and elaborated once, never created
        in a \code{generate} loop whose bounds are not constant;
  \item \kw{modport} on every port, so that directions are explicit;
  \item no tasks with timing controls (\code{@}, \code{\#}) in anything that
        is synthesised, and no clocking blocks in the design half;
  \item no dynamic types, no classes, no covergroups in an interface that
        reaches synthesis; put those in a separate verification interface
        bound to the design (\cref{sec:beyond:bind});
  \item parameters, packed structs and enums are fine.
\end{itemize}

\begin{ruleofthumb}
Use interfaces freely at the testbench boundary and between blocks you own.
Think twice before putting one on the top-level pins of a chip or on a
deliverable IP block.
\end{ruleofthumb}

\subsection{Compared with \VHDL{} records and mode views}

\VHDL{}-2008 lets you put a record in a port, which solves half the problem:

\begin{rosetta}
\begin{rleft}
type bus_req_t is record
  req   : std_logic;
  addr  : std_logic_vector(11 downto 0);
  wdata : std_logic_vector(31 downto 0);
end record;

entity ram is
  port (clk   : in  std_logic;
        req_i : in  bus_req_t;
        rsp_o : out bus_rsp_t);
end entity;
\end{rleft}
\begin{rright}
module ram (simple_bus.slave b);
  // b.req, b.addr, b.wdata,
  // b.gnt, b.rdata ...
endmodule
\end{rright}
\end{rosetta}

The \VHDL{} record port bundles signals, but every element of a record port
has the mode of the port itself. A bidirectional bus therefore needs two
record types, one per direction, and you are back to naming things twice. The
record also cannot carry the write procedure.

\VHDL{}-2019 addresses this directly with \emph{interfaces} and \emph{mode
views}\index{VHDL-2019}, which are close in spirit to a modport: a view
declaration gives each record element its own direction, and a port declared
with that view behaves like a modport. It is the right feature. The problem
is support: at the time of writing, mode views are implemented in only a
handful of tools, and no mainstream synthesis flow accepts them. Treat
\VHDL{}-2019 as a language you can read but not yet ship.

% ===========================================================================
\section{Clocking blocks}
\label{sec:beyond:clocking}
\index{clocking block}

Every \VHDL{} testbench you have written contains a line like
\code{wait for 2 ns;} placed just after a clock edge, so that the stimulus
does not change at the same instant the design samples it. That convention is
the entire race-avoidance strategy of a classical \VHDL{} testbench, and it is
carried in the author's head, not in the language.

\SV{} makes it a declaration. A \kw{clocking} block names a clock, lists the
signals that are read and written with respect to that clock, and attaches a
\emph{skew} to each direction. It is declared inside an interface, which is
why it appears in a design chapter at all.

\begin{svcode}[caption={A clocking block fixes sampling and driving skew},label={lst:beyond:cb}]
interface simple_bus (input logic clk, input logic rst_n);
  logic req, gnt;
  logic [11:0] addr;

  clocking cb @(posedge clk);
    default input #1step output #1ns;
    output req, addr;
    input  gnt;
  endclocking

  modport tb (clocking cb, input clk, rst_n);
endinterface
\end{svcode}

One line declares both skews. \code{input \#1step} samples every signal listed
as \code{input} one \emph{time step} before the clock edge, so a reader sees
the value the design saw and never the value the design has just produced.
\code{output \#1ns} delays every assignment to a clocking-block output until
\SI{1}{\nano\second} after the edge, in the middle of the cycle, where no
flip-flop can be confused by it. Using the block is unremarkable:
\code{@(bus.cb)} waits for the clocking event, \code{bus.cb.addr <= 12'h040}
drives, and \code{bus.cb.gnt} reads.

Note the direction convention, because it trips up everyone once. A clocking
block is written from the point of view of the \emph{testbench}, so what the
testbench drives is \code{output}, even though those same wires are inputs to
the design.

That is as much as a design reader needs. A clocking block is not
synthesisable, and the rest of the subject belongs to the testbench:
\refverificationbook explains what \code{\#1step} means in terms of the scheduler
regions, why the drive skew is what stops the stimulus racing the design, and
what happens when a clocking-block output and a plain procedural assignment
target the same signal. The point to take away here is structural. The timing
relationship between testbench and design is now written down once, inside the
interface that carries the protocol, instead of being repeated as a hand-tuned
delay in every stimulus block.

% ===========================================================================
\section{Declaring intent: \texttt{always\_comb}, \texttt{always\_ff}, \texttt{always\_latch}}
\label{sec:beyond:always}
\index{always\_comb}\index{always\_ff}\index{always\_latch}

\VHDL{} has one process. \SV{} has one general \code{always} plus three
specialised forms whose only purpose is to let you state what you meant so
that the tool can check it. This is a language feature with no \VHDL{}
counterpart: in \VHDL{} an accidental latch is caught by a synthesis warning,
if you read your logs.

\begin{svcode}
always_comb begin              // combinational; tool checks it
  y = a & b;
end

always_ff @(posedge clk or negedge rst_n) begin
  if (!rst_n) q <= '0;
  else        q <= d;
end

always_latch begin
  if (en) q <= d;              // latch, and you said so on purpose
end
\end{svcode}

What is actually checked differs between the three, and the differences
matter.

\code{always\_comb} has an inferred sensitivity list that includes everything
read in the block \emph{and everything read inside any function it calls}.
That last part is the fix for the classic \code{always @(a or b)} bug where a
signal was left out of the list. It also runs once at time zero regardless of
any event, so the outputs are consistent from the start; a \VHDL{} process
with a sensitivity list does the same, so this will feel natural. Two further
restrictions: no timing controls are allowed inside it, and the variables it
writes may not be written by any other process. That second rule is stronger
than anything in \VHDL{}, and it turns a whole class of multiple-driver bugs
into a compile error.

\code{always\_ff} requires exactly one event control and the same
single-writer rule. What the \LRM{} does \emph{not} require is that you use
non-blocking assignments inside it; that is a lint rule, not a language rule.
Tools warn, and you should treat the warning as an error.

\code{always\_latch} checks nothing you could not check yourself. Its value is
documentary: it tells the next reader that the latch is intentional, and it
lets a lint rule flag every latch that is \emph{not} inside one.

\begin{ruleofthumb}
In new \SV{} \RTL{}, plain \code{always} appears only in testbenches. Every
synthesisable process is \code{always\_comb}, \code{always\_ff} or, rarely,
\code{always\_latch}.
\end{ruleofthumb}

\subsection{\texttt{unique}, \texttt{unique0} and \texttt{priority}}
\index{unique}\index{priority}

The same idea applies to decisions. \VHDL{} case statements must be
exhaustive, which the compiler enforces; \SV{} case statements need not be,
so the language offers modifiers that state your assumption instead.

\begin{svcode}
unique case (sel)          // exactly one branch matches, no overlap
  2'b00: y = a;
  2'b01: y = b;
  2'b10: y = c;
  2'b11: y = d;
endcase

priority if (irq_hi)  grant = HI;   // at least one matches;
else if (irq_mid) grant = MID;      // first match wins
else if (irq_lo)  grant = LO;
\end{svcode}

\code{unique} asserts two things: at least one branch matches, and no two
branches match simultaneously. \code{unique0} drops the first requirement, so
"none of them" is legal but overlap is not. \code{priority} asserts only that
at least one branch matches, and keeps the written order.

At simulation time these become runtime checks: the simulator issues a
violation report when the assumption fails, with the values that caused it. At
synthesis time they become directives, the modern replacement for the
\code{// synopsys parallel\_case} and \code{full\_case} pragmas. That is where
the danger lies.

\begin{gotcha}{\texttt{unique} is a promise, not a check, in synthesis}
Synthesis believes you. If you write \code{unique case} over a
one-hot-encoded signal and the signal is momentarily not one-hot in silicon,
the synthesised logic may produce something that the \RTL{} simulation never
produced, because synthesis used your promise to optimise away the
multiplexer priority. The simulator flags the violation; the gates do not.
Only write \code{unique} when a separate mechanism makes the claim true: an
assertion, a formal proof, or a structurally one-hot encoding.
\end{gotcha}

While on the subject of case: \SV{} inherits \code{casex} and \code{casez}
from Verilog, and adds \code{case ... inside}, which uses wildcard matching in
one direction only.

\begin{gotcha}{\texttt{casex} treats an X in the \emph{expression} as a wildcard}
\code{casez} treats \code{Z} (and \code{?}) in the case \emph{items} as
don't-care, which is what you want. \code{casex} additionally treats \code{X}
and \code{Z} in the \emph{selecting expression} as don't-care. So an
uninitialised or corrupt signal carrying \code{X} matches the first branch and
simulation carries on as though nothing were wrong, hiding exactly the bug
that \code{X} was there to reveal. \VHDL{}'s \code{std\_logic} case statement
would have forced you to write a \code{when others} and think about it. Never
write \code{casex}. Prefer \code{case (x) inside} with \code{4'b1???} style
items, which is both readable and X-safe.
\end{gotcha}

% ===========================================================================
\section{Packed structures and unions}
\label{sec:beyond:struct}
\index{struct!packed}\index{union!packed}

This is the construct \VHDL{} designers miss most once they have used it, and
it is the one to adopt first.

A \kw{packed} struct is a named bit vector. The fields occupy contiguous bits,
the first-declared field is the most significant, and the whole object behaves
as a single vector wherever a vector is expected. It can cross a port, sit in
an array, be assigned from a literal, be cast to and from a plain vector, and
be indexed by name.

\begin{svcode}[caption={An instruction word decoded by field name},label={lst:beyond:instr}]
typedef struct packed {
  logic [6:0] funct7;
  logic [4:0] rs2;
  logic [4:0] rs1;
  logic [2:0] funct3;
  logic [4:0] rd;
  logic [6:0] opcode;
} instr_r_t;                    // exactly 32 bits, MSB first

module decoder (
  input  logic [31:0] iword,
  output logic [4:0]  rd_o,
  output logic        is_add_o
);
  instr_r_t ins;
  assign ins       = instr_r_t'(iword);   // a cast, not a conversion
  assign rd_o      = ins.rd;
  assign is_add_o  = (ins.opcode == 7'b0110011) &&
                     (ins.funct3 == 3'b000) && (ins.funct7 == '0);
endmodule
\end{svcode}

Compare this with the same decoder written against bit indices. The
\SV{} version cannot silently shift by one when someone widens \code{funct7},
and \code{\$bits(instr\_r\_t)} tells you the width without counting. The cast
\code{instr\_r\_t'(iword)} is free: it reinterprets bits, it does not convert
values, and it is a compile error if the widths differ.

Packed structs nest, cross ports, and form arrays:

\begin{svcode}
typedef struct packed {
  logic        valid;
  logic [3:0]  id;
  instr_r_t    ins;
} pipe_stage_t;                  // 1 + 4 + 32 = 37 bits

pipe_stage_t stage [0:3];        // unpacked array of packed structs
pipe_stage_t [3:0] pipe;         // packed array: one 148-bit object

always_ff @(posedge clk) begin
  stage[0] <= '{valid: 1'b1, id: next_id, ins: instr_r_t'(iword)};
  for (int i = 1; i < 4; i++) stage[i] <= stage[i-1];
end
\end{svcode}

The \code{'\{...\}} form is an assignment pattern, and naming the fields makes
it order-independent. \code{'\{default: '0\}} zeroes an arbitrarily
complicated structure in one line, which is worth knowing for reset branches.

\begin{rosetta}
\begin{rleft}
type instr_r_t is record
  funct7 : std_logic_vector(6 downto 0);
  rs2    : std_logic_vector(4 downto 0);
  rs1    : std_logic_vector(4 downto 0);
  funct3 : std_logic_vector(2 downto 0);
  rd     : std_logic_vector(4 downto 0);
  opcode : std_logic_vector(6 downto 0);
end record;

-- and a hand-written function
-- to_record(slv) : 20 lines of
-- slicing, maintained by hand.
\end{rleft}
\begin{rright}
typedef struct packed {
  logic [6:0] funct7;
  logic [4:0] rs2;
  logic [4:0] rs1;
  logic [2:0] funct3;
  logic [4:0] rd;
  logic [6:0] opcode;
} instr_r_t;

// and the conversion is
instr_r_t'(iword)
\end{rright}
\end{rosetta}

The \VHDL{} record is a fine data structure, but it has no defined bit layout,
so moving between a record and a \code{std\_logic\_vector} means writing and
maintaining a pair of conversion functions. That is the whole difference:
\kw{packed} gives the record a layout.

\begin{gotcha}{Unpacked structs are not vectors}
Leave out the word \kw{packed} and you get an unpacked struct. It is still
useful in testbenches, where fields can be queues and strings, but it has
no bit layout, cannot be cast to a vector, cannot cross a synthesisable port
as a vector, and cannot be assigned from a literal vector. Forgetting
\kw{packed} on an \RTL{} type produces a wall of type errors that never
mention the missing word.
\end{gotcha}

\subsection{Packed unions}

A packed union gives two names to the same bits:

\begin{svcode}
typedef union packed {
  instr_r_t     r;              // decoded as an R-type
  instr_i_t     i;              // decoded as an I-type
  logic [31:0]  raw;            // or as a plain word
} instr_t;

instr_t w;
assign w.raw = iword;
assign rs1   = w.r.rs1;         // same bits, different view
\end{svcode}

Every member of a packed union must have the same width; the compiler
enforces it. This is the clean way to express a register that has a
field-oriented view and a byte-oriented view, or an instruction word with
several formats.

\SV{} also has \emph{tagged} unions, which add a hidden tag field and let the
compiler check that you read the member you last wrote, with
\code{case ... matches} for discrimination. They are elegant and almost
unsupported: several major simulators accept them and most synthesis tools do
not. Know they exist; do not build a design on them.

% ===========================================================================
\section{Enumerations that carry a base type}
\label{sec:beyond:enum}
\index{enumeration}

A \VHDL{} enumeration is an abstract type whose encoding is chosen by
synthesis unless you attach an attribute. A \SV{} enumeration is a set of
named constants of a base type you choose, and the encoding is part of the
declaration:

\begin{svcode}
typedef enum logic [2:0] {
  IDLE = 3'b001, RUN = 3'b010, DONE = 3'b100
} state_e;                       // explicitly one-hot

typedef enum logic [1:0] { A, B, C } small_e;  // 0, 1, 2 by default
\end{svcode}

The base type also fixes what happens on reset: \code{logic}-based
enumerations start at \code{X}, which is what you want, because an
uninitialised state machine should be visibly broken. Using \kw{bit} as the
base type would start it at zero and hide the missing reset.

Enumerations have built-in methods, which have no \VHDL{} counterpart at all
beyond \code{'image} and \code{'succ}:

\begin{svcode}
state_e s = RUN;
$display("%s", s.name());   // "RUN"  -- the reason to use them in $display
s = s.next();               // DONE   (wraps round after the last)
s = s.prev();               // RUN
$display("%0d", s.num());   // 3      -- number of members
state_e f = s.first();      // IDLE
state_e l = s.last();       // DONE
\end{svcode}

\code{.name()} alone justifies the construct: every waveform dump, every
log message and every assertion failure can print the state by name with no
lookup table. \code{.next()} and \code{.prev()} wrap around, so a loop written
as "advance until you reach \code{last()}" runs forever if you write the exit
test wrongly; iterate with \code{for (int i = 0; i < s.num(); i++)} instead.

\begin{notebox}{Enumerations are strongly typed, up to a point}
You cannot assign an integer to an enumeration variable without an explicit
cast: \code{s = state\_e'(2)} compiles, \code{s = 2} does not. But the cast is
not checked: \code{state\_e'(3'b111)} produces a variable holding a value
with no name, and \code{.name()} returns an empty string for it. This is
weaker than \VHDL{}, where the type simply has no such value.
\end{notebox}

% ===========================================================================
\section{Parameterised types and elaboration-time computation}
\label{sec:beyond:params}
\index{parameter!type}

\VHDL{} generics carry values. \SV{} parameters carry values \emph{and types}:

\begin{svcode}[caption={A FIFO parameterised by the type it carries},label={lst:beyond:paramtype}]
module fifo #(
  parameter type T     = logic [7:0],
  parameter int  DEPTH = 8
)(
  input  logic clk, rst_n,
  input  logic push, pop,
  input  T     din,
  output T     dout,
  output logic full, empty
);
  T mem [DEPTH];
  localparam int PTRW = $clog2(DEPTH);
  // ...
endmodule

// instantiated with a struct as the payload:
fifo #(.T(pipe_stage_t), .DEPTH(16)) u_f ( ... );
\end{svcode}

The FIFO now works for any packed type without knowing its width, and
\code{\$bits(T)} recovers the width where it is needed. In \VHDL{} the same
effect requires a generic package (2008) or a generic subprogram (2019), both
of which are more verbose and, more importantly, only partially supported:
generic packages work in the major commercial simulators, \tool{QuestaSim}
among them, but synthesis support has historically been patchy, and generic
subprograms are 2019 and effectively unavailable.

Related tools:

\begin{itemize}
  \item \code{type(expr)}\index{type() operator} yields the type of an
        expression, so \code{type(din) tmp;} declares a variable that follows
        the port. It is useful in macros and in generic code.
  \item \code{\$bits(x)}\index{\$bits} gives the width in bits of a type or an
        expression, including a whole struct.
  \item Functions are evaluated at elaboration when their arguments are
        constant, so you can compute parameters:
\end{itemize}

\begin{svcode}[caption={A constant function initialising a localparam array},label={lst:beyond:constfn}]
function automatic int unsigned atan_rom_entry(int unsigned i);
  // evaluated at elaboration time; no hardware is built
  return int'($atan(2.0**(-real'(i))) * (2.0**13));
endfunction

localparam int unsigned ATAN [0:13] =
  '{ atan_rom_entry(0),  atan_rom_entry(1),  atan_rom_entry(2),
     atan_rom_entry(3),  atan_rom_entry(4),  atan_rom_entry(5),
     atan_rom_entry(6),  atan_rom_entry(7),  atan_rom_entry(8),
     atan_rom_entry(9),  atan_rom_entry(10), atan_rom_entry(11),
     atan_rom_entry(12), atan_rom_entry(13) };
\end{svcode}

This is the same technique \code{cordic\_pkg} in \refexamplebook uses to build
its arctangent table, and it is one place where \SV{} and \VHDL{} are genuinely
equivalent: a \VHDL{} constant initialised by a pure function does the same job,
and does it with better type checking.

The real package differs from the sketch above in three details, and each of
them matters. It scales by $2^{15}$, not $2^{13}$, because the rotator's
internal angle format \code{iangle\_t} carries two guard bits below the
external Q2.13 format. It rounds with \code{\$rtoi(x + 0.5)} rather than
truncating with \code{int'(x)}, which is worth up to one LSB on every entry.
And it writes a \code{for} loop instead of fourteen calls. The first difference
alone scales the whole table by four: the sketch above starts 6433, 3798,
2006, \dots{} while \code{ATAN\_ROM} starts 25736, 15193, 8027, \dots{} Read
the real one before you copy either.

\begin{tipbox}{Write the loop, not the list}
A \code{for} loop inside a constant function returning an unpacked array is
cleaner than the fourteen calls above, and every current simulator and
synthesis tool supports it. Write \code{localparam int ATAN [0:13] =
build\_atan();} and put the loop in \code{build\_atan}. That is what
\code{cordic\_pkg} does. \tool{Icarus Verilog} is the exception: it cannot
return an unpacked array from a function, which is why \refsimulationbook shows
how to generate the same table from Python instead.
\end{tipbox}

% ===========================================================================
\section{\texttt{let}, \texttt{`define}, and the preprocessor}
\label{sec:beyond:macro}
\index{`define}\index{let}

Parameters and constant functions cover almost everything a design needs to
compute before elaboration. What they do not cover is text, and for text
\SV{} has a C-style preprocessor. \VHDL{} has nothing of the kind before 2019,
and you have not missed much.

\code{`define} performs textual substitution with no type, no scope and no
respect for the hierarchy. A macro defined in one file is visible in every
file compiled afterwards in the same compilation unit, which means two IP
blocks that both define \code{`WIDTH} will fight. Use \code{parameter} and
\code{localparam} for values, \code{typedef} for types, and reserve
\code{`define} for things a parameter genuinely cannot express: include
guards, tool-conditional code, and macros that must expand to declarations.

\code{let} covers many of the remaining cases with proper scoping and proper
argument semantics:

\begin{svcode}
let handshake(v, r) = v && r;          // scoped, typed, no textual games
let is_pow2(n)      = (n & (n-1)) == 0;

always_comb accept = handshake(req_valid_i, req_ready_o);
\end{svcode}

A \code{let} is expanded at each use with the arguments bound at that point,
so it can refer to signals in the scope where it is used, unlike a function.
It is legal in assertions, which is where it earns its keep.

\begin{gotcha}{Macros are text, and text does not respect precedence}
\code{`define SUM(a,b) a + b} followed by \code{x = `SUM(1,2) * 3;} yields
\code{1 + 2 * 3}, that is \code{7}. Always parenthesise both the arguments and
the whole body: \code{`define SUM(a,b) ((a) + (b))}. This is the oldest bug in
C and it arrived in \SV{} unchanged.
\end{gotcha}

% ===========================================================================
\section{\texttt{generate} and \texttt{bind}}
\label{sec:beyond:bind}
\index{generate}\index{bind}

Everything so far has changed what you can write \emph{inside} a module. The
last two constructs change what you can build \emph{out of} modules, and the
second of them is the most useful thing in this chapter.

\code{generate} needs little explanation, because you already have it:
\code{for}, \code{if} and \code{case} generate blocks map one-to-one onto the
\VHDL{} constructs of \cref{sec:rosetta:generate}. Two differences are worth
remembering. The loop variable is declared \kw{genvar} rather than being
implicit in the loop, and naming the block, \code{: g\_lane} after the loop
header, is strongly advised because the name becomes part of the hierarchical
path. An unnamed generate block gets a tool-invented name such as
\code{genblk3}, and every waveform path, every \code{bind} target and every
constraint that mentions it then depends on how many unnamed blocks happen to
precede it in the file.

\kw{bind} is the interesting one, and it has no \VHDL{} equivalent whatsoever.

\begin{svcode}[caption={Binding a checker to a design without touching its source},label={lst:beyond:bindex}]
module handshake_chk (input logic clk, rst_n,
                      input logic valid, ready);
  property no_drop;
    @(posedge clk) disable iff (!rst_n)
      valid && !ready |=> valid;
  endproperty
  a_no_drop: assert property (no_drop)
    else $error("valid dropped before ready");
endmodule

// Attach it to every instance of the module type:
bind cordic_rot handshake_chk u_chk_req (
  .clk(clk_i), .rst_n(rst_ni),
  .valid(req_valid_i), .ready(req_ready_o));

// ...or to one specific instance only:
bind u_dut.u_cordic handshake_chk u_chk_rsp (
  .clk(clk_i), .rst_n(rst_ni),
  .valid(rsp_valid_o), .ready(rsp_ready_i));
\end{svcode}

\kw{bind} instantiates a module (or an interface, or a program) \emph{inside}
another module's scope, from outside, without editing that module. Inside the
bound instance, the target's internal signals are visible by name, so a
checker can watch a state register that never reaches a port.

Why this matters: it is the mechanism that makes assertion-based verification
possible on \RTL{} you do not own. You receive a block from another team, or
from a vendor, and you attach a hundred assertions and a functional coverage
model to it in a separate file that your flow can strip out for synthesis. In
\VHDL{} your only options are to edit the source or to be content with what
the ports expose.

\begin{tipbox}{Keep bound checkers in their own file list}
Put \code{bind} statements and checker modules in a \file{*\_bind.sv} file
compiled only for simulation and formal. The design file list stays clean, the
synthesis tool never sees an assertion, and nobody has to remember a
\code{`ifndef SYNTHESIS} guard.
\end{tipbox}

\begin{gotcha}{A bound port connection is elaborated in the target's scope}
The expressions in a \code{bind} port list are resolved inside the target
module, not where the \code{bind} statement is written. That is the whole
point, and it is also the trap: a name that happens to exist in the target
with a different width is connected and silently truncated or extended,
exactly as an ordinary port connection would be, and a name that does not
exist in the target is an elaboration error pointing at a file that contains
no such signal. Connect every port by name, and check the widths by hand the
first time.
\end{gotcha}

Two things in \cref{lst:beyond:bindex} are not explained here. The
\kw{property} and the \code{assert property} belong to SystemVerilog
Assertions, which is \refverificationbook; and what a bound checker is actually
worth, in a testbench that is already running constrained-random stimulus, is
the subject of \refverificationbook. This chapter owns the construct: \code{bind}
attaches a module to a module. The next two chapters spend it.

% ===========================================================================
\section{Checklist}
\label{sec:beyond:checklist}
% ===========================================================================

\begin{itemize}
  \item \textbf{Use an interface} for any bundle of more than about four
        signals that crosses more than one module boundary. Give it a
        \kw{modport} per role.
  \item \textbf{Use a packed struct} wherever you would have written a
        \VHDL{} record on a port. It crosses a port, it casts to a vector,
        and the field names survive into the waveform.
  \item \textbf{Give every enumeration a base type.} The tool will not
        choose one for you the way a \VHDL{} synthesiser does.
  \item \textbf{\code{always\_comb}, \code{always\_ff}, \code{always\_latch}},
        never a bare \code{always}, in \RTL{}. They are free assertions.
  \item \textbf{\code{unique} and \code{priority} are promises}, and the
        synthesiser believes them. Only write one where you can prove it.
  \item \textbf{Compute constants at elaboration} with a constant function
        rather than pasting a table into the source, and check that your
        tools accept it before you rely on it.
  \item \textbf{Declare the testbench timing, do not time it by hand.} Where
        a \VHDL{} testbench writes \code{wait for 2 ns;} after every edge, an
        \SV{} interface carries a clocking block with
        \code{default input \#1step output \#1ns} and the question does not
        come up again. \refverificationbook uses it.
  \item \textbf{Name every generate block.} \code{: g\_lane} after the loop
        header, always. The name becomes the hierarchical path, and
        \code{genblk3} is not a path you want in a script.
  \item \textbf{\code{bind} a checker} onto somebody else{\rq}s \RTL{} rather
        than editing it. There is no \VHDL{} equivalent and it is one of the
        most useful things in the language.
  \item \textbf{Prefer a parameter to a \tick{define}.} The preprocessor has
        no scope and no types.
\end{itemize}
